\section{\texorpdfstring{Almost $X$-uniform distributions support $X$-uniform distributions}{Almost X-uniform distributions support X-uniform distributions}}
\label{sec:almost_support_uniform}

We prove in this section \expref{Theorem}{thm:almost_support_uniform},
which states that
if $\mu$ is an almost $X \times X$-uniform distribution with
small enough error, then there exists an $X$-uniform distribution $\mu'$
supported on the support of $\mu$. We recall the statement of the theorem.

\restate{\expref{Theorem}{thm:almost_support_uniform}}{
Let $\mu$ be a distribution supported on a set $S \subset G$ which is
almost $(X \times X)$-uniform with error $\eps < 0.5 |X|^{-7}$. Then
there exists a distribution $\mu'$ supported on $S$ which is $X$-uniform.
}

Fix such a distribution $\mu$, and let $S$ denote its support, $S=\{g:
\mu(g)>0\}$. Let $R_X$ be the representation of $G$ acting on $X$. By
\expref{Proposition}{prop:uniform-by-convex-hull}, $S$ supports an $X$-uniform
distribution iff $R_X(U_G)=\E_{g \in G}[R_X(g)]$ lies in the convex
hull of $\{R_X(g): g \in S\}$. Assume this is not the case; then there
exists 
%
a hyperplane 
$H$ which passes through $R_X(U_G)$ and such that
all $\{R_X(g): g \in S\}$ lie on one side of $H$.

We first project $H$ into 
%
a hyperplane 
with a simpler representation. Let
$R_X \equiv e_0 \one + e_1 R_1 + \cdots + e_t R_t$ denote the
decomposition of $R_X$ into unitary nonequivalent irreducible
representation, and let $d_i = \dim(R_i)$ denote the dimension of each irreducible representation.
Essentially, we will project $H$ to ``use'' only one copy from each nontrivial
irreducible representation. That is, we will show that $H$ can be
projected to a hyperplane separating $0$ from $\{R_1(g) \times \cdots
\times R_t(g): g \in S\}$.

\begin{proposition}\label{prop:separating_zero_R}
There exists a map $L:G \to \R$ given by
$$
L(g) := \sum_{i \in [t]} \sum_{j,k \in [d_i]}
\lambda_{i,j,k} \cdot R_i(g)_{j,k}
$$
for some coefficients $\{\lambda_{i,j,k}
\in \C: i \in [t], j,k \in [d_i]\}$ such that
\begin{enumerate}
\item $\E_{g \in G}[L(g)]=0$;
\item for all $g \in S$, $L(g)>0$.
\end{enumerate}
\end{proposition}

\begin{proof}
The existence of a hyperplane $H$ separating $R_X(U_G)$ from
$\{R_X(g): g \in S\}$ implies that there exists a map $L':G \to \R$ defined as
$L'(g) = \sum_{x,y \in X} \alpha_{x,y} R_X(g)_{x,y}$
for some real coefficients $\{\alpha_{x,y} \in \R: x,y \in X\}$ such that
$$
L'(g) > \E_{h \in G}[L'(h)]
$$
for all $g \in S$. Applying the linear transformation mapping $R_X$
into the basis of irreducible representations, we get that $L'(g)$
can be expressed as
$$
L'(g) = \sum_{\ell \in [e_0]} \beta_{0,\ell} \one(g)
+ \sum_{i \in [t]} \sum_{j,k \in [d_i]} \sum_{\ell \in [e_i]}
\beta_{i,j,k,\ell} R_i(g)_{j,k}\,,
$$
where $\beta_{0,\ell},\beta_{i,j,k,\ell} \in \C$ are obtained by a linear
transformation (over $\C$) of $\alpha_{x,y}$. Observe that $\E_{g \in
G}[L'(g)] = \sum_{\ell \in [e_0]} \beta_{0,\ell}$ by Schur's lemma,
and define $L(g) := L'(g) - \E[L'(g)]$. Note that $L:G \to \R$ is real
since $L':G \to \R$ was real, that $\E[L]=0$ and that $L(g)>0$ for all
$g \in S$. The coefficient $\lambda_{i,j,k}$ is given by the sum of all
$\beta_{i,j,k,\ell}$ over $\ell \in [e_i]$.
\end{proof}

We may assume without loss of generality that $\E_{g \in G}[L^2(g)]=1$ by multiplying all
coefficients $\lambda_{i,j,k}$ by an appropriate factor. The main idea
is to show that if $\mu$ is almost $X \times X$ uniform, then $\E_{g
\sim \mu}[L^2(g)] \approx \E_{g \in G}[L^2(g)]=1$ while $\E_{g \sim
\mu}[L(g)] \approx \E_{g \in G}[L(g)]=0$. Combining this with a bound
on $\|L\|_{\infty}$ a simple calculation shows that it cannot be the
case that $L(g)>0$ for all $g$ in the support of $\mu$.

The first step is to show that the
coefficients $\lambda_{i,j,k}$ cannot be very large.

\begin{proposition}\label{prop:bound_coefficients_lambda}\hspace{1mm}
\[
\sum_{i \in [t]}\sum_{j,k \in [d_i]} \frac{|\lambda_{i,j,k}|^2}{d_i} = 1\,.
\]
In particular, $|\lambda_{i,j,k}| \le |X|^{1/2}$ for all $i,j,k$.
\end{proposition}

\begin{proof}
We assumed $\E[L^2]=1$, which, since $L$ is
real, implies $\E[|L|^2]=1$. Using Schur's lemma we get
\begin{eqnarray*}
1 &=& \E_{g \in G}\bigl[L(g) \cdot \overline{L(g)}\bigr] \\
&=& \sum_{i,i' \in [t]} \sum_{j,k \in [d_i]} \sum_{j',k' \in [d_{i'}]}
\lambda_{i,j,k} \overline{\lambda_{i',j',k'}} \E_{g \in G}\bigl[R_i(g)_{j,k}
\overline{R_{i'}(g)}_{j',k'}\bigr]\\
& =& \sum_{i \in [t]} \sum_{j,k \in [d_i]} \frac{|\lambda_{i,j,k}|^2}{d_i}\,,
\end{eqnarray*}
and in particular $|\lambda_{i,j,k}|^2 \le d_i \le |X|$.
\end{proof}

An immediate corollary is that $L(g)$ can never be very large.

\begin{corollary}\label{cor:L_bound_infty}
$|L(g)| \le |X|^{2.5}$ for all $g \in G$.
\end{corollary}

\begin{proof}
We have $|R_i(g)_{j,k}| \le 1$ since $R_i$ is unitary, hence $|L(g)|
\le \sum_{i \in [t]} \sum_{j,k \in [d_i]} |\lambda_{i,j,k}|
\le |X|^{2.5}$ since $\sum d_i^2 \le |X|^2$.
\end{proof}

The bound on $|\lambda_{i,j,k}|$ together with the assumption
that $\mu$ is almost $X \times X$-uniform, implies that the first
and second moments of $L$ are approximately
the same under $\mu$ and under the uniform distribution over $G$.

\begin{proposition}\label{prop:L_moments_approx}
Let $\mu$ be a distribution which is
almost $X \times X$-uniform with error $\eps$. Then
\begin{enumerate}
\item $|\E_{g \sim \mu}[L(g)]| \le \eps |X|^{4.5}$;
\item $|\E_{g \sim \mu}[L^2(g)]-1| \le \eps |X|^7$.
\end{enumerate}
\end{proposition}

\begin{proof}
We have
$$
|\E_{g \sim \mu} [L(g)]| \le \sum_{i \in [t]} \sum_{j,k \in [d_i]}
|\lambda_{i,j,k}| |\E_{g \sim \mu} [R_i(g)_{j,k}]|\,.
$$
The bound on the first moment follows since $\sum d_i^2
\le |X|^2$; since $|\lambda_{i,j,k}| \le |X|^{1/2}$ by
\expref{Proposition}{prop:bound_coefficients_lambda}; and since $\mu$
is in particular $X$-uniform with error $\eps |X|$, we have by
\expref{Proposition}{prop:rephrase_conditions_dist_irr} that $|\E_{g \sim
\mu}[R_i(g)_{j,k}]| \le \eps |X|^2$. The bound on the second moment is
proved in a similar way. Recall that we proved the identity $\sum |\lambda_{i,j,k}|^2/d_i=1$ in
\expref{Proposition}{prop:bound_coefficients_lambda}. So
\begin{eqnarray*}
\E_{g \sim \mu} [|L(g)|^2-1] &=& \sum_{i,j,k}
|\lambda_{i,j,k}|^2 \cdot \left(\E_{g \sim \mu}[|R_i(g)_{j,k}|^2] - 1/d_i\right)\\
&+& \sum_{(i,j,k) \ne (i',j',k')} \lambda_{i,j,k}
\overline{\lambda_{i',j',k'}} \cdot
\E_{g \sim \mu}[R_i(g)_{j,k} \overline{R_{i'}(g)_{j',k'}}]\,.
\end{eqnarray*}
To conclude the proof we need to show that $\E_{g \sim
\mu}[|R_i(g)_{j,k}|^2] \approx 1/d_i$ and that $\E_{g \sim
\mu}[R_i(g)_{j,k} R_{i'}(g)_{j',k'}] \approx 0$.

The condition that $\mu$ is
almost $X \times X$-uniform with error $\eps$ is equivalent to
$$
\|R_{X \times X}(\mu) - R_{X \times X}(U_G)\|_{\infty} \le \eps\,.
$$
Switching to the Frobenius norm, this implies
$$
\|R_{X \times X}(\mu) - R_{X \times X}(U_G)\|_{\Fr} \le \eps |X|^2\,.
$$
We now decompose $R_{X \times X}$ to simpler representations, coming
from the irreducible representations of $R_X$. We have that $R_{X
\times X}=R_X \otimes R_X$, which since $R_X$ is real, also gives $R_{X
\times X} = R_X \otimes \overline{R_X}$.  Now, if $R_X \equiv e_0 \one
+ \sum_{i=1}^t e_i R_i$ is the decomposition of $R_X$ into irreducible
unitary nonequivalent representations, we have $$ R_{X \times X} \equiv
e_0^2 \one + \sum_{i=1}^t e_0 e_i (R_i+\overline{R_i}) + \sum_{i,i'=1}^{t}
e_i e_{i'} (R_{i} \otimes \overline{R_{i'}})\,.$$ Note that this is not
the decomposition of $R_{X \times X}$ into irreducible representations,
since $R_i \otimes \overline{R_{i'}}$ is reducible! Nevertheless, we
observe that as the basis change for $R_X$ was unitary, so is the basis
change for $R_{X \times X}$ (since the tensor product of two unitary
matrices is again unitary). In particular, we get that 
\[
\|(R_i \otimes
\overline{R_{i'}})(\mu) - (R_i \otimes \overline{R_{i'}})(U_G)\|_{\Fr}
\le \eps |X|^2\,,
\]
which implies that
$$
\|(R_i \otimes \overline{R_{i'}})(\mu) -
(R_i \otimes \overline{R_{i'}})(U_G)\|_{\infty} \le \eps |X|^2\,.
$$
The matrix $R_i \otimes \overline{R_{i'}}$ is
indexed by $((j,j'),(k,k'))$, where $(R_i \otimes
\overline{R_{i'}})(g)_{(j,j'),(k,k')} = R_i(g)_{j,k}
\overline{R_{i'}(g)_{j',k'}}$. We thus get that for any $i,j,k,i',j',k'$
we have
$$
\left| \E_{g \sim \mu}[R_i(g)_{j,k} \overline{R_{i'}(g)_{j',k'}}]
- \E_{g \in G}[R_i(g)_{j,k} \overline{R_{i'}(g)_{j',k'}}]
\right| \le \eps |X|^2\,.
$$
To conclude, by Schur's lemma
$$
\E_{g \in G}[R_i(g)_{j,k} \overline{R_{i'}(g)_{j',k'}}]
= \frac{1}{d_i} 1_{(i,j,k)=(i',j',k')}\,.
$$
The bound for the second moment now follows by
elementary calculations analogous to the ones for the first moment.
\end{proof}

We are now ready to prove \expref{Theorem}{thm:almost_support_uniform}.

\begin{proof}[Proof of \expref{Theorem}{thm:almost_support_uniform}]
Let $\mu$ be almost $X \times X$ uniform
with error $\eps \le 0.5 |X|^{-7}$.
Summarizing \expref{Corollary}{cor:L_bound_infty} and
\expref{Proposition}{prop:L_moments_approx}, we have
\begin{enumerate}
\item $\|L\|_{\infty} \le |X|^{2.5}$;
\item $E_{g \sim \mu}[L(g)] \le \eps |X|^{4.5}$;
\item $E_{g \sim \mu}[L(g)^2] \ge 1 - \eps |X|^{7}$.
\end{enumerate}
However, since we assumed by contradiction
that $L(g)>0$ for all $g$ in the support of $\mu$, we have
$$
\E_{g \sim \mu}[L(g)^2] \le \|L(g)\|_{\infty}
\cdot \E_{g \sim \mu}[L(g)] \le |X|^{2.5} \cdot \eps |X|^{4.5},
$$
\ie, we have
$$
1 - \eps |X|^7 \le \eps |X|^7,
$$
which is false whenever $\eps < 0.5 |X|^{-7}$.
\end{proof}
